% titel: Blatt 6, A6: Polynome 3-ten Grades (Cardano)
% Teil II: der durchgerechnete Weg.

\subsection*{Gegeben und gesucht}
Für $p,q\in\mathbb R$ mit $D=q^2+p^3>0$ und $u=-q+\sqrt D$, $v=-q-\sqrt D$:
nachweisen, dass $x=\sqrt[3]{u}+\sqrt[3]{v}$ (reelle Kubikwurzel)
Nullstelle von $z^3+3pz+2q$ ist; danach alle Nullstellen von
$z^3+3\sqrt[3]{9}z+8$ bestimmen.

\subsection*{a) Nachweis, dass $x$ Nullstelle ist}
\emph{Gedankenschritt:} Der Schlüssel ist die Identität
$uv=-p^3$. Denn damit wird das Binom $(a+b)^3=a^3+b^3+3ab(a+b)$ so
zurückgeführt, dass der $3px$-Term entsteht.
Zuerst:
\begin{equation*}
  uv=(-q+\sqrt D)(-q-\sqrt D)=q^2-D=q^2-(q^2+p^3)=-p^3,
\end{equation*}
also (mit der reellen Kubikwurzel) $\sqrt[3]{uv}=\sqrt[3]{-p^3}=-p$. Mit
$a=\sqrt[3]{u}$, $b=\sqrt[3]{v}$ und $x=a+b$ gilt $ab=\sqrt[3]{uv}=-p$;
das Binom liefert
\begin{align*}
  x^3&=a^3+b^3+3ab\,(a+b)=u+v+3\underbrace{\sqrt[3]{uv}}_{\mathstrut -p}\,x\\
     &=(-2q)+3(-p)x=-2q-3px .
\end{align*}
Umgestellt: $x^3+3px+2q=0$, also ist $x$ eine Nullstelle.

\subsection*{b) Nullstellen von $z^3+3\sqrt[3]{9}\,z+8$}
\emph{Gedankenschritt:} $p$ und $q$ aus der Standardform ablesen:
$3p=3\sqrt[3]{9}$, also $p=\sqrt[3]{9}$; $2q=8$, also $q=4$.
\begin{equation*}
  D=q^2+p^3=16+(\sqrt[3]{9})^3=16+9=25,\qquad \sqrt D=5,
\end{equation*}
\begin{equation*}
  u=-4+5=1,\qquad v=-4-5=-9,
  \qquad x_1=\sqrt[3]{1}+\sqrt[3]{-9}=1-\sqrt[3]{9}.
\end{equation*}
Faktorisieren mit $x_1$:
\begin{equation*}
  z^3+3\sqrt[3]{9}\,z+8
  =(z-x_1)\bigl(z^2+x_1z+(x_1^2+3\sqrt[3]{9})\bigr).
\end{equation*}
Die quadratische Formel für die restlichen Nullstellen:
\begin{equation*}
  z_{2,3}=\frac{-x_1\pm\sqrt{x_1^2-4(x_1^2+3\sqrt[3]{9})}}{2}
  =\frac{-x_1\pm i\sqrt{3x_1^2+12\sqrt[3]{9}}}{2}.
\end{equation*}
Mit $x_1=1-\sqrt[3]{9}$ vereinfacht sich
\begin{equation*}
  3x_1^2+12\sqrt[3]{9}
  =3\bigl(1+2\sqrt[3]{9}+\sqrt[3]{9}^2\bigr)
  =3\bigl(1+\sqrt[3]{9}\bigr)^2,
\end{equation*}
also
\begin{equation*}
  z_{2,3}=\frac{\sqrt[3]{9}-1\pm i\sqrt3\,(1+\sqrt[3]{9})}{2}.
\end{equation*}

\begin{ergebnis}
\begin{equation*}
  z_1=1-\sqrt[3]{9},\qquad
  z_{2,3}=\frac{\sqrt[3]{9}-1\,\pm\,i\sqrt3\,(1+\sqrt[3]{9})}{2}.
\end{equation*}
\end{ergebnis}

\begin{kontrolle}
\textbf{Numerisch:} $z_1\approx-1{,}0801$, $z_{2,3}\approx0{,}540\pm
2{,}667i$. Einsetzen bestätigt alle drei Nullstellen (siehe Verifikation:
$|p(z)|$ ist für alle drei verschwindend klein). \checkmark

\textbf{Vieta-Probe:} Das Polynom hat keinen $z^2$-Term, also ist die
Summe aller Nullstellen $0$; ferner summiert $\sum z_iz_j=3\sqrt[3]{9}$.
Beide Passen zur gefundenen Lösung. \checkmark

\textbf{Grenzfall:} $D=25>0$, also genau eine reelle plus zwei konjugiert
komplexe Nullstellen -- konsistent mit der Cardano-Voraussetzung.
\checkmark
\end{kontrolle}

\begin{merksatz}
Beim Cardano-Verfahren zuerst $\sqrt[3]{uv}=-p$ berechnen; damit wird
$x^3=-2q+3\sqrt[3]{uv}\,x=-2q-3px$ ein kurzer Beweis der
Nullstelleneigenschaft. Faktorisieren liefert die restlichen Nullstellen
über die quadratische Formel. Für $D>0$: eine reelle, zwei komplexe
Nullstellen.
\end{merksatz}
